%% Currículo atualizado — Renan Aragão
%% Versão voltada a recrutadores de posições Full-Stack, Laravel, Vue e IA aplicada.
\documentclass[10pt,a4paper]{article}
\usepackage[utf8]{inputenc}
\usepackage[T1]{fontenc}
\usepackage{lmodern}
\usepackage[margin=1.5cm,top=1.2cm,bottom=1.15cm]{geometry}
\usepackage{xcolor}
\usepackage{titlesec}
\usepackage{enumitem}
\usepackage{tabularx}
\usepackage[hidelinks]{hyperref}

\definecolor{accent}{HTML}{D95020}
\definecolor{ink}{HTML}{17202B}
\definecolor{muted}{HTML}{566574}
\renewcommand{\familydefault}{\sfdefault}
\pagestyle{empty}
\setlength{\parindent}{0pt}
\setlength{\parskip}{3pt}

\titleformat{\section}{\large\bfseries\color{ink}}{}{0em}{}[{\color{accent}\titlerule[0.7pt]}]
\titlespacing*{\section}{0pt}{8pt}{3pt}
\setlist[itemize]{leftmargin=1.15em,itemsep=1pt,topsep=2pt,parsep=0pt}
\newcommand{\job}[3]{\textbf{#1} \hfill {\color{muted}\small #3}\\{\color{muted}#2}\par\vspace{-2pt}}
\newcommand{\link}[2]{\href{#1}{\color{accent}#2}}

\begin{document}
\fontsize{9.4}{11.2}\selectfont

{\Huge\bfseries\color{ink} Renan Aragão}\\[3pt]
{\large\color{accent}\bfseries Desenvolvedor Full-Stack · Laravel, Vue, microserviços e IA aplicada}\\[3pt]
{\small\color{muted} Fortaleza, CE · oportunidades presenciais, híbridas ou remotas}\\[2pt]
{\small \link{https://renanaragao9.github.io/renanaragao9}{Portfólio} \quad
\link{https://linkedin.com/in/renan-aragao/}{LinkedIn} \quad
\link{https://github.com/renanaragao9}{GitHub} \quad
\link{https://seuracha.aragaolabs.com.br}{SeuRacha} \quad
\link{https://focimor.aragaolabs.com.br}{Focimor}}

\section{Perfil}
Desenvolvedor full-stack com mais de 3 anos de experiência em produtos web, integrações e automações de negócio. Trabalho principalmente com PHP/Laravel no back-end, Vue/Nuxt no front-end e Python para OCR, modelos de linguagem e automações. Valorizo escopo e critério de pronto claros, alinhamento técnico antes da implementação e entregas testáveis.

\section{Destaques}
\begin{itemize}
  \item Participei da criação de um agente em Python que lê, extrai e classifica documentos com Tesseract OCR, spaCy e Transformers, poupando centenas de horas da equipe de auditoria.
  \item Atuei em um chatbot com LLM conectado à base de dados da empresa, tornando a consulta a informações de projetos mais rápida.
  \item Automatizei a geração de boletos para cerca de \textbf{15 mil associados}, reduzindo erros e trabalho manual do financeiro.
  \item Participei do desenvolvimento de \textbf{mais de seis sistemas web} para o setor de energia.
  \item Em uma colaboração freelancer anterior, participei do desenvolvimento de um sistema de ponto e RH utilizado por mais de 30 empresas.
\end{itemize}

\section{Experiência}
\job{Desenvolvedor Full-Stack · 3e Soluções}{Fortaleza, CE · híbrido · energia e eficiência energética}{nov/2024 -- atual}
\begin{itemize}
  \item Participação no desenvolvimento de sistemas web com Laravel, Vue 3 e APIs integradas para métricas, análise de dados e gestão de custos.
  \item Soluções em Python para OCR, classificação de documentos e chatbot com LLM; uso de FastAPI, HuggingFace e OpenRouter.
  \item Microsserviços e integrações com Kafka, Kong e Docker; deploy com Laravel Vapor e AWS (S3, CloudFront e SES).
  \item Atuação com MySQL, MongoDB e Firebase em aplicações web e mobile.
\end{itemize}

\job{Desenvolvedor Full-Stack · CABEMCE}{Fortaleza, CE · instituição com mais de 15 mil associados}{ago/2023 -- out/2024}
\begin{itemize}
  \item Participei do desenvolvimento de funcionalidades para ERP interno, incluindo cadastro de dependentes, histórico de assistência e permissões.
  \item Automação em Python para geração de boletos e apoio ao faturamento de cerca de 15 mil associados.
  \item Integrações REST com OAuth/JWT, PostgreSQL e rotinas de atendimento e assistência.
  \item Trabalho em equipe com Kanban e Gitflow no ciclo de entrega.
\end{itemize}

\job{Estagiário de TI · Tribunal Regional do Trabalho da 7ª Região}{Fortaleza, CE}{jun/2022 -- ago/2023}
\begin{itemize}
  \item Suporte N1 a N3, consultas e diagnósticos em Oracle e análise de causa-raiz para reduzir chamados indevidos direcionados ao desenvolvimento.
\end{itemize}

\section{Projetos em produção}
\begin{itemize}
  \item \textbf{\link{https://seuracha.aragaolabs.com.br}{SeuRacha}} \textit{(projeto pessoal)}: plataforma de esporte amador com rachas, reserva de quadras, torneios e gestão de arenas. \emph{Laravel 13, Filament, Next.js e TypeScript.}
  \item \textbf{\link{https://focimor.aragaolabs.com.br}{Focimor}} \textit{(projeto pessoal)}: plataforma de adoção de pets com match entre pessoas e ONGs, chat e painel administrativo. \emph{Laravel 13, Filament, Next.js e TypeScript.}
\end{itemize}

\section{Tecnologias}
\begin{tabularx}{\linewidth}{@{}l X@{}}
\textbf{Back-end} & PHP 8, Laravel (Vapor, Queues, Nova, Filament, Sanctum, Passport), APIs REST \\
\textbf{Front-end} & Vue 3, Pinia, Nuxt, Next.js, TypeScript, Tailwind CSS \\
\textbf{Python e IA} & Python, FastAPI, Flask, Tesseract OCR, spaCy, Transformers, RAG, LLMs, HuggingFace \\
\textbf{Dados e infra} & PostgreSQL, MySQL, MongoDB, Firebase, Redis, Oracle, Docker, AWS, Kafka, Kong \\
\textbf{Qualidade} & PHPUnit, pytest, Git/Gitflow, Kanban \\
\end{tabularx}

\section{Formação}
\textbf{Bacharelado em Ciência da Computação} · Faculdade Multiversa, Fortaleza/CE \hfill {\color{muted}\small 2020 -- 2024}

\end{document}
